Nullniveau der Lageenergie: der tiefste Punkt des Tals. Zustand 1: auf dem ersten Hügel (nur Lageenergie). Zustand 2: zuunterst im Tal (nur kinetische Energie). Zustand 3: auf dem zweiten Hügel (Lageenergie und kinetische Energie). Nur die Höhen zählen, nicht die Form der Bahn.

\begin{abcliste}
\abc
Sei $H$ die Höhe des ersten Hügels und $v_1$ die Geschwindigkeit zuunterst. Die ganze Lageenergie im Zustand 1 ist im Zustand 2 zu kinetischer Energie geworden:
\begin{align}
 m\,g\,H &= \frac{1}{2}\,m\,v_1^2 \\
 v_1 &= \vaF \\
   &= \sqrt{2\cdot\g\cdot\Hh} \\
   &= \va \approx \result{\vaP{0}{2}}
\end{align}

\abc
Sei $h_2$ die Höhe des zweiten Hügels. Der Wagen hat nur einen Teil seiner Höhe verloren; dieser Teil ist zu kinetischer Energie geworden. Zwischen den Zuständen 1 und 3:
\begin{align}
 m\,g\,H &= m\,g\,h_2 + \frac{1}{2}\,m\,v_2^2 \\
 v_2 &= \vbF \\
   &= \sqrt{2\cdot\g\cdot(\Hh-\hz)} \\
   &= \vb \approx \result{\vbP{0}{2}}
\end{align}
\end{abcliste}
